\hypertarget{link_to_finj}{}{\Large \bf f\_inj}

{\it Computes $f$, an estimate of the frequency (in per unit) at a given bus, from the evolution of the rectangular components $v_x$ and $v_y$ of the bus voltage. A measurement time constant $T$ is involved. Can be used in the model of an injector only.}

~

Syntax : \hspace*{10ex}\parbox[t]{\textwidth}{\sf 
\& f\_inj \\
name of variable $f$\\
data name, parameter name or math expression for $T$
}

~ \hrule ~

Internal states : $v_{xm}$ ~ and ~ $v_{ym}$.

Discrete variables : none

Equations :
\begin{eqnarray}
\dot{v}_{xm} & = & \frac{v_x - v_{xm}}{T} \label{vxmdot} \\
\dot{v}_{ym} & = & \frac{v_y - v_{ym}}{T} \label{vymdot} \\
0 & = & \omega_{ref, pu} + \frac{(v_y - v_{ym})v_{xm} - (v_x - v_{xm}) v_{ym}}{2 \pi f_{N} T (v_{xm}^2 + v_{ym}^2)} - f \label{f}
\end{eqnarray}

Initialization of internal states : $v_{xm} = v_{x}$ ~ and ~ $v_{ym}=v_y$.

~

{\bf Explanation of model}

The model uses $v_x$ and $v_y$ as inputs but these variables are automatically inherited and must not be declared. Only the name given to the estimated frequency and the time constant $T$ must be provided.

As shown in the figure below, $v_x$ and $v_y$ are the projections of the bus voltage phasor on the references axes $x$ and $y$ rotating at the angular speed $\omega_{ref}$ (rad/s). The latter is known from the settings of the simulation. The angular frequency of the corresponding rotating vector is given by :
\[ \omega = \omega_{ref} + \frac{d\phi}{dt} \]
with:
\[ \phi = \arctan \frac{v_y}{v_x} \]

\begin{center}
\includegraphics{codegen/phasor.pdf}
\end{center}

The frequency in per unit is given by :
\[ f = \frac{\omega}{2 \pi f_N} = \frac{\omega_{ref}}{2 \pi f_N} + \frac{1}{2 \pi f_N} \frac{d\phi}{dt}   
= \omega_{ref, pu} + \frac{1}{2 \pi f_N} \frac{d}{dt} \left( \arctan \frac{v_y}{v_x} \right) \]
where $f_N$ is the nominal frequency (known from the system data).

By developing the last term, it is easily found that:
\begin{equation}\label{f2}
f = \omega_{ref, pu} + \frac{1}{2 \pi f_N} \frac{\dot{v_y} v_x - \dot{v_x} v_y}{v_x^2 + v_y^2}
\end{equation}

To filter the transients affecting $v_x$ and $v_y$, a measurement device with a time constant $T$ is simulated. The ``measured'' values, denoted $v_{xm}$ and $v_{ym}$, evolve according to (\ref{vxmdot}, \ref{vymdot}). These measured values and their derivatives are then used in (\ref{f2}), which becomes:
\[ f = \omega_{ref, pu} + \frac{1}{2 \pi f_N} \frac{\dot{v}_{ym} v_{xm} - \dot{v}_{xm} v_{ym}}{v_{xm}^2 + v_{ym}^2} \]
Replacing $\dot{v}_{xm}$ and $\dot{v}_{ym}$ by their expressions (\ref{vxmdot},\ref{vymdot}) yields Eq. (\ref{f}).

\vspace*{2ex}
A recommended value for $T$ is in the order to $0.05 - 0.10$ s. A zero value for $T$ is not allowed. If too small a value is specified for $T$, the solver may encounter a singularity and the simulation may not proceed.

